\chapter{Mechanising Lilac Part II: Proofs}
\emph{Having overcome the fundamental change in the definitions, consolidate inconsistencies and instantiating Iris, we present here the 2nd aspect
of a mechanisation project: A LOT of theorem proving. We present the soundness proofs of Lilac proof rules drawing attention to 3 of them
in particular: \texttt{wp\_frame}, \texttt{wp\_bind} and \texttt{wp\_unif}, due to each of them influencing the definition of the \texttt{wp}
connective (in order to validate their soundness). \texttt{wp\_bind} motivated the use Mathlib's Probability Kernels, with kernel style definitions
and proofs. \texttt{wp\_unif} the central and most complex proof in the entire project, motivates the creation of an entire sub-library of
measure theoretic results about measurable spaces obtained through finite footprint on the Hilbert Cube. Some of the development of these proofs
happened simultaneously with surprising progress in LLM-based auto-formalisation in Lean. We report on how these tools were used with success in simpler
sub proofs (which are time consuming to work out but not conceptually hard), but they are not particularly good or successful, when the rich grounding of Mathlib's
Api is not present and a library of lemmas needs to first be built in order to carry out a proof as was the case for \texttt{wp\_unif}. It is worth noting
that as past contributor of Mathlib, Mathlib has exacting code quality standards, for the generalisability and suitable abstractions used in contributed
code, and as seen the success of such autoformalisation is in part due to the availability of such an excellent library. The ambitions of this mechanisation
is to forma a foundation for further developments in probabilistic separation logic, such extended Lilac (with conditioning) and Bayesian Separation Logic \cite{},
so we try to maintain good abstractions and remark on the extensibility at points in the chapter.}

\Todo{Proof irrelavance for closed\_krm\_subtype also for conditioning}
\Todo{Pi-lambda theorem not suitable for BaSL}

Let's recall in one place how we now define variables, contexts the language and the logic, by going through a simple but important proof rule \texttt{wp\_ret}.
The syntax and semantics of the language has not changed and is still expressed as a measurable function $\sem{\Delta} \marrow A$, where
$ \sem{\Delta} = A \otimes \ldots \otimes A$, where A represents types endowed with a measurable space. That's the semantics of a program
while in the syntax to refer to variable we use dependent de brujin indices.

The logic which originally started out in the same manner as the language, i.e. a deep embedding with dependent de brujin indices, has transformed into
a shallow embedding with higher order syntax style treatment of the variables. Variables are of the type \texttt{RV A}, which expands to $HC \marrow A$,
with the standard borel measurable space on the domain (hilbert cube) and additionally the finite footprint property, whhich we state now for the first time:

% The aspect of playing with definitions sneaks into this chapter as well,
% with some modifications to the \texttt{wp} connective. The complexity of definition of the \texttt{wp} connective is motivated by establishing the soundness
% of the above the 3 rules}



\section{To overview}
\subsection{Reviewing Probabilistic Separation Logics}
More recently there has been a lot of research into probabilistic separation logics: \citep{og_psl}, introduced separation as probabilistic independence.
\citep{lilac} introduced Lilac, which allowed usage of continuous distributions, providing a measure theoretic definition of resources. \citep{basl} introduced reasoning about bayesian
updates (supporting the \texttt{observe} and \texttt{normalize} constructs), thereby providing an axiomatic semantics for the general form of first order
probabilistic programs (ie, BPPL as introduced by \citep{prob-semantics}).

There is a parallel line of work in relational separation logics. One could argue that relational logics are more important in the probabilistic domain, since in bayesian
probabilistic programing (so including \texttt{observe} and \texttt{normalize}), frequently generates probability distributions which have no closed form solution! Sampling
from these can only be done using approximation methods like Variational Inference. The justification for relational logics classical is that relational properties are
expressed more naturally, but could perhaps still be done using unary reasoning, in an inelegant way. In probabilistic programs however, relational reasoning might
in fact be more expressive making it important. Relational reasoning for probabilistic programs (though not using a separation logic) can be found here \cite{relational-logic}.
More recently there has been a novel combination of relational and unary reasoning for probabilistic programs in a separation logic by \citep{bluebell}.

\section{Pi-lambda-theorem as an induction principle}
\label{pi-lambda}
The $\pi$-$\lambda$-theorem is tool in measure theory which can be used for equality of two measures. We use it prove a certain uniqueness of measures theorem,
used to establish that separating conjunction yeilds a PCM (partial commutative monoid) structure.
It is similar to an induction principle on the structure of a $\sigma$-algebra. So we need prove some property C holds in the base and inductive cases
of a $\sigma$-algebra. However there isn't an inductive type underlying this
induction principle! The reason is two fold 1) when generating a measurable set from some base measurable sets, a specific
set can be shown to be measurable by multiple different permutations of application of the axioms. For reference below
is included an inductive definition generating a $\sigma$-algebra from base measurable seats.

\begin{leancode}
/-- The smallest σ-algebra containing a collection `s` of basic sets -/
inductive GenerateMeasurable (s : Set (Set α)) : Set α → Prop
  | protected basic : ∀ u ∈ s, GenerateMeasurable s u
  | protected empty : GenerateMeasurable s ∅
  | protected compl : ∀ t, GenerateMeasurable s t → GenerateMeasurable s tᶜ
  | protected iUnion : ∀ f : ℕ → Set α, (∀ n, GenerateMeasurable s (f n)) →
      GenerateMeasurable s (⋃ i, f i)
\end{leancode}

2) The $\pi$-$\lambda$-theorem provides a relaxation on our proof obligation. We dont prove that the property C holds for countable unions 
(\texttt{iUnion} above) if it holds for each of the sets we're taking a union over. Instead we only need to do so for countably pairwise disjoint unions.

Being closed under disjoint unions instead of unions, is the difference between a $\lambda$-system compared to a $\sigma$-algebra. This relaxtion is the 
content of the $\pi$-$\lambda$-theorem which we don't go into here.
However there is an alternative "inductive" definition of a $\lambda$ system with the following axioms: 
\begin{enumerate}
\item The system $\mathcal{C}$ containt the base measurable sets we chose
\item be closed under proper difference, for $\forall V, U \in \mathcal{C}, V \subseteq U \rightarrow U \setminus V \in \mathcal{C}$
\item closed under increasing limits, take $i \in \mathbb{N}$ and $U_i \in \mathcal{C}$ then $\bigcup_{i \in \mathbb{N}} U_i \in \mathcal{C}$
\end{enumerate}

Unfortunately this is not the definition of $\lambda$-system in Mathlib, and it doesn't provide the "induction principle" according to this definition.
For the case of probability measures ($\mu$ (universal set) = 1), there is a simple workaround, which is sufficient for Lilac. However BaSL uses unnormalised
measures since it support bayesian conditioning, and so, there may be no simple workaround but reproving this alternative induction principle.

% \section{Binder Discipline}
% The choice between Dependent De Brujin Indices and Parametric Abstract Higher Order Syntax (PHOAS).
% \subsection{Dependent De Brujin Indices}
% \begin{itemize}
% \item More straightforward to the uninitiated following a little more intuitively from the paper. Important if
% we want this is to be an accessible code base for future probabilistic separation logics to follow from.
% \item Measurability of maps expressible explicitly
% \end{itemize}

% \subsection{PHOAS}
% \begin{itemize}
% \item Notion of substitution inbuilt
% \item Example programs written in it are naturally readable. No separate parser and translation from a normal
% lambda calculus required. Correctness of the parser doesn't need to be proven... (but further consideration on
% this point requires deciding if there will be a separate executable implementation of the APPL syntax)
% \end{itemize}


% \section{Equality of dependently typed records in Lean}


% \section{Choice of Kripke Resource Monoid (KRM) over Unital ordered Resource algebra (ORA) }
% ORA's are established as the standard by former Rocq-iris projects, but I couldn't see why not just directly
% use the KRM. It was conceptual simplicity traded for some possibly unweildy proofs. The unweildiness turned out to be
% restricted to a single file?

% \section{Deep vs shallow embedding}

% \section{Support for recursion over the data structure even with (shallow?/deep? see SSProve talk) embedding}
% (compare with SProp usage in Rocq for ORA? or something...)
% Even though the monotonicity property did not require induction, substitution atleast as defined in the paper requires it.

% The aim, is to do things as given in the paper, and then think about if substitution could be done away with? If removing
% the ability to do recursion over assertions is desirable, etc

% \section{Strong abstractions across the proof development}
% \subsection{The language}
% This is built on the existing abstractions. We use the typeclass \texttt{DenotationalMeas} to decouple the syntax from
% the language from the logic. The language only needs to me able to interpreted as measurable maps with a few additional
% requirements as given by \texttt{DenotationalMeas}.

% \subsection{Resource monoid}
% \emph{for the version of Lilac with/without conditioning}

% \subsection{Separate treatment of BI connectives and Probability specific connectives}

% \section{Metaprogramming and tactical support}

% - The correct type for meaurable spaces because type casting issue, it needed to be a dependent type.

% - Correctly instantiating BI, modelling random variables, "intrinsically type calculus"? 
% - Talk about Lilac, give the intuition that this is the actual line by line semantics